\documentclass[11pt,reqno]{amsart}
\usepackage[T1]{fontenc}
\usepackage{lmodern,amsmath,amssymb,amsthm,mathtools,microtype}
\usepackage[margin=1.1in]{geometry}
\usepackage{xr-hyper}
\usepackage[hidelinks]{hyperref}
\externaldocument[TC-]{two_collision}[two_collision.pdf]
\allowdisplaybreaks[2]
\setlength{\emergencystretch}{2em}
\numberwithin{equation}{section}
\newtheorem{theorem}{Theorem}[section]
\newtheorem{lemma}[theorem]{Lemma}
\newtheorem{proposition}[theorem]{Proposition}
\newtheorem{corollary}[theorem]{Corollary}
\theoremstyle{definition}\newtheorem{definition}[theorem]{Definition}
\theoremstyle{remark}\newtheorem{remark}[theorem]{Remark}
\newcommand{\R}{\mathbb R}\newcommand{\Z}{\mathbb Z}
\newcommand{\E}{\mathbb E}\newcommand{\dd}{\,\mathrm d}
\newcommand{\one}{\mathbf1}\newcommand{\csch}{\operatorname{csch}}
\newcommand{\arcosh}{\operatorname{arcosh}}\newcommand{\tr}{\operatorname{tr}}
\newcommand{\diag}{\operatorname{diag}}\newcommand{\supp}{\operatorname{supp}}

\title[Geometric collision thresholds]{Geometric stability and curvature identification\\for uniform collision thresholds}
\author{Qian Qi}
\date{September 9, 2026}
\subjclass[2020]{37D50, 37A50, 60G55, 70H09}
\keywords{Lorentz gas, collision thresholds, specular trajectories, marked counting laws,
 hyperbolic periodic orbit, parameter response}
\begin{document}
\begin{abstract}
We prove a collision-order-uniform threshold law for periodic dispersing
billiards in an open neighborhood of circular scatterers. The complete
maximal-count event splits into the finitely many competing shortest
channels. A channel with gap $g$ and endpoint curvatures $\kappa_0,\kappa_1$
contributes $\varepsilon_+^2/(2A\sinh(j\gamma))$ times a smooth relative
amplitude, where $\cosh\gamma=\sqrt{(1+g\kappa_0)(1+g\kappa_1)}$.
The collar and every fixed normalized remainder bound are independent of
the number $j$ of complete flights. The proof combines complete-event
localization, an alternating-curvature Dirichlet reduction, and a relative
trace-norm determinant estimate. It permits unequal curvatures, nonsymmetric
boundary jets, and changes of the minimizing pair. Smooth additive marks
have a differentiated expansion with an exact axial source-convergence
domain. In an explicit three-parameter noncircular family all threshold
locations are identical, while their unlabelled amplitudes determine the
contact-curvature multiset and the free area by an absolutely convergent
M\"obius inversion. We specify complete-record topology and prove response
for smooth record tests and transverse moving sampling cuts. The circular specialization gives explicit instability ratios and a
quadratic onset distinct from that of independent flights.
\end{abstract}
\maketitle
\section{Geometric thresholds and curvature data}
\label{sec:geometric-results}
The location of a shortest-flight threshold is metric information. Its
probability also records the transverse instability of the reflected
segment. These data coincide in a one-parameter circular family, since
the gap already determines the radius. They separate under deformations
of the scatterer. We prove a threshold law on an open set of noncircular
geometries, including the competition between distinct shortest pairs,
and show that its amplitudes determine curvature information which all
threshold locations leave undetermined.

Fix the triangular lattice $\Lambda$ and a reference radius
$R_*\in(0,1/2)$. A scatterer $D_\xi$ is described by its support function
$h_\xi(\theta)$, where $\xi$ ranges over a compact subset $K_0$ of a
finite-dimensional parameter domain. The function $(\xi,\theta)\mapsto
h_\xi(\theta)$ is $C^\infty$ on a neighborhood of $K_0\times S^1$.
There is a $C^4$ neighborhood of the constant $R_*$ in which the
following conclusions hold. No reflection or rotational symmetry is
assumed. Uniform constants for a prescribed derivative order may depend
on the corresponding higher norms of this fixed smooth family; they are
not uniform over unbounded higher derivatives in a $C^4$ ball.

Write $\mathcal E=\{e\in\Lambda:|e|=1\}$, so $|\mathcal E|=6$.
The unique closest segment from $D_\xi$ to $D_\xi+e$ has length
$g_e(\xi)$ and endpoint curvatures $\kappa_{e,0}(\xi)$ and
$\kappa_{e,1}(\xi)$. All these functions are smooth. Put
\begin{equation}\label{eq:g-general-data}
 \begin{split}
 A(\xi)&=\sqrt3/2-|D_\xi|,\qquad g_*(\xi)=\min_{e\in\mathcal E}g_e(\xi),\\
 c_{e,a}&=1+g_e\kappa_{e,a}\quad(a=0,1),\qquad
 c_e=\sqrt{c_{e,0}c_{e,1}},\qquad \gamma_e=\arcosh c_e.
 \end{split}
\end{equation}
Reversing $e$ exchanges the curvatures and leaves $g_e,c_e,\gamma_e$
unchanged. The full preparation is normalized unit-speed phase volume
on $(\R^2/\Lambda)\setminus D_\xi$. We write $N_t$ for the actual
collision count under that preparation.

\begin{theorem}[Stability and competition of collision thresholds]
\label{thm:g-stability}
For all such compact smooth families there is a collar
$\varepsilon_0>0$, independent of every integer $j\ge1$, such that,
if $0\le t-jg_*<\varepsilon_0$, then
\begin{equation}\label{eq:g-competition}
 \begin{split}
 \Pr_\xi(N_t\ge j+1)&=\Pr_\xi(N_t=j+1)\\
 &=\sum_{e\in\mathcal E}
   \frac{d_{j,e,+}^{\,2}}{2A(\xi)\sinh(j\gamma_e(\xi))}
   \bigl[1+d_{j,e,+}H_{j,e}(\xi,d_{j,e,+})\bigr],\\
 d_{j,e}&=t-jg_e(\xi).
 \end{split}
\end{equation}
Here $d_+=\max(d,0)$, and every fixed mixed derivative of $H_{j,e}$
in $\xi$ and its nonnegative offset is bounded independently of $j$.
The functions have smooth right extensions at offset zero. For
$t\le jg_*$ the probability is zero. The collar is smaller than
$\min_{K_0}g_*/4$.

More precisely, the $e$ summand is the probability of the oriented
alternating channel. At its own window $t=jg_e+d$, $0<d<\varepsilon_0$,
that channel has the same formula, whether or not $e$ minimizes $g_e$.
In this last statement the channel specifies the facing contact arcs
of all $j+1$ impacts, not the entire count event at that window.
\end{theorem}

The formula uses the six smooth lengths separately, not derivatives of
the possibly nonsmooth minimum $g_*$. It therefore retains changes of
the minimizing pair. It also keeps the exponentially small coefficient
relative to its own error. For a circle, $c_e=(1-R)/R$ and
$1/(2A)=\lambda_R/(4R)$; summing its six equal terms gives
Theorem~\ref{thm:unweighted} without a change of preparation or horizon.
The complete circular proof is given below, followed by the geometric
proof in Sections~\ref{sec:g-localization}--\ref{sec:g-integration}.

Let $x_{e,0}(\xi),x_{e,1}(\xi)$ be the normal outgoing states, and
let $f_\xi$ be a smooth additive collision mark in their neighborhoods.
Set $a_{e,b}=f_\xi(x_{e,b})$, $\bar f_e=(a_{e,0}+a_{e,1})/2$, and
$w_{j,e}=\sum_{i=0}^j a_{e,i\bmod2}$.

\begin{theorem}[Geometric marked response and its exact source domain]
\label{thm:g-marked}
At the channel's own threshold the marked probability is
\begin{equation}\label{eq:g-marked}
 Z_{j,e}(\xi,q,d)=
 \frac{d^2e^{q w_{j,e}(\xi)}}{2A(\xi)\sinh(j\gamma_e(\xi))}
       [1+dH_{j,e}^{f}(\xi,q,d)].
\end{equation}
Every fixed mixed derivative of the normalized remainder is bounded
uniformly in $j$ on compact parameter and source sets. The complete
marked maximal-count event in the collar of Theorem~\ref{thm:g-stability}
is the sum of these expressions at the positive offsets $d_{j,e}$.
For fixed $\xi$ and real $q$, and for any sufficiently small fixed
$d\ge0$, the series of normalized channel probabilities
$\sum_{j,e}Z_{j,e}(\xi,q,d)/d^2$ converges if and only if
\begin{equation}\label{eq:g-source-domain}
               \max_{e\in\mathcal E}\{q\bar f_e(\xi)-\gamma_e(\xi)\}<0.
\end{equation}
At $d=0$ use right limits. On compact subsets of this open source
domain every fixed mixed derivative can be summed term by term.
The same assertion holds for finitely many real sources, replacing
$q\bar f_e$ by their scalar product.
\end{theorem}

The series concerns channel thresholds in distinct windows. It is not
an unrestricted long-time pressure. Its exact source boundary is useful
precisely because the contributing trajectories have been identified.
The conservative circular source condition of Theorem~\ref{thm:marked}
remains valid and is a sufficient subset of \eqref{eq:g-source-domain}.

Our inverse consequence is not a new interpretation of the radius in a
circle. In Section~\ref{sec:g-inverse} we construct a three-parameter
family of noncircular scatterers with the same shortest gaps at all six
orientations. Their unlabelled count amplitudes determine the unordered
triple of contact curvatures, with the free area also recovered rather
than assumed. In the symmetric one-parameter subfamily, just two
amplitudes suffice. Section~\ref{sec:g-record-response} specifies a
supremum metric for complete records and proves differentiated estimates
for smooth tests and a transverse moving cut. Those are distinct from
convergence for bounded discontinuous decoders.

\paragraph{Organization and notation.}
Sections~\ref{sec:results}--\ref{sec:consequences} treat the circular
specialization, with radius $R$ and the notation introduced there.
Sections~\ref{sec:g-localization}--\ref{sec:g-record-response} prove
the geometric statements for $D_\xi$. The circular convention
$g_R=1-2R$ is not imposed on the latter family. The final section
compares the two threshold regimes with other dynamical limit and
inverse problems.

\input{sections/01_results}
\section{Stationary flux and the complete marked event}
\label{sec:flux}
For an outgoing boundary state write $\dd s=R\dd\alpha$ and
$p=\sin\phi$. Its unnormalized flux is $\dd s\dd p$, of total mass
$4\pi R$. A flight tube has phase volume $\dd s\dd p\dd r$.
Except for a null set of directions, a line on the lattice torus meets
the open disk in both time directions: irrational lines are dense, and
the remaining rational directions are countable. Tangencies and their
finite regular preimages have zero flux. The positive gap prevents
finite-time accumulation. Thus the flight tubes cover phase volume up
to a null set, also when corridors give infinite horizon. It follows that
\begin{equation}\label{eq:santalo}
 \overline\tau_R=\int\tau_R\dd\nu=\frac{A_R}{2R},\qquad
 \dd\mu_R=\lambda_R\dd\nu(x)\dd s,\quad 0\le s<\tau_R(x).
\end{equation}
Reflection preserves flux, and $B_R$ is invertible and $\nu$-preserving
on the full-measure regular section.

For completeness, integration of the allowable suspension positions gives
\begin{equation}\label{eq:tails-new}
 \Pr_R(N_t\ge k)=\lambda_R\int
       \bigl[(t-S_{k-1}\tau_R)_+-(t-S_k\tau_R)_+\bigr]\dd\nu.
\end{equation}
Indeed, from $(x,s)$ the $k$th hit occurs at $S_k\tau_R(x)-s$.
The length of the allowed part of $[0,\tau_R(x))$ is
$(t-S_{k-1}\tau_R(B_Rx))_+-(t-S_k\tau_R(x))_+$; invariance proves
\eqref{eq:tails-new}. This is a stationary suspension identity, not a
renewal theorem. Since $S_k\tau_R\ge kg_R$, only finitely many tails
are nonzero at each fixed time. Their sum telescopes to
$\E_RN_t=\lambda_Rt$.

A marked version must not shift the first term by invariance while leaving
its mark unchanged. We use the following direct form instead.

\begin{proposition}[First-impact coordinates]\label{prop:marked-flux}
Let $t=jg_R+\varepsilon$, where $j\ge1$ and
$0<\varepsilon<g_R$. Let $y$ be the outgoing state at the first future
impact and let $r$ be its time. For any bounded Borel function $G$ of the
complete impact record,
\begin{equation}\label{eq:record-flux}
 \E_R[\one_{\{N_t=j+1\}}G]
 =\lambda_R\int_{S_j\tau_R(y)<t}
       \int_0^{t-S_j\tau_R(y)}
       G\bigl(\mathcal P_{j,R}(y,r;t)\bigr)\dd r\dd\nu(y).
\end{equation}
Here $\mathcal P_{j,R}$ is the actual path record, with impact times
$r+S_i\tau_R(y)$, $0\le i\le j$, and collision states $B_R^iy$.
The formula includes the initial and terminal free-flight portions.
\end{proposition}
\begin{proof}
Replace the last collision state $x$ in \eqref{eq:santalo} by $y=B_Rx$
and the suspension height by the first residual time
$r=\tau_R(x)-s$. The measure is $\lambda_R\dd\nu(y)\dd r$ with
$0<r<\tau_R(B_R^{-1}y)$. The event of at least $j+1$ hits is
$r+S_j\tau_R(y)\le t$. Its upper bound on $r$ is at most
$\varepsilon<g_R\le\tau_R(B_R^{-1}y)$, so no additional truncation
occurs. A $(j+2)$nd hit would require $t\ge(j+1)g_R$, which is
impossible. This proves the integration limits and the event identity.

Write $q_i,v_i^+,v_i^-$ for the impact position and the outgoing and
incoming velocities. Starting at $q_0-rv_0^-$, the initial segment lasts
$r$. Subsequent straight segments and specular reflections are prescribed
by $B_R^iy$. After the last impact, the path continues for
$t-r-S_j\tau_R(y)\ge0$. Both end portions are shorter than $g_R$ and
cannot contain an additional impact. This gives the claimed complete
mechanical record, including both velocity traces at each hit.
\end{proof}

Taking $G=1$ recovers
\begin{equation}\label{eq:Aj}
 P_j(R,\varepsilon):=\Pr_R(N_{jg_R+\varepsilon}=j+1)
   =\lambda_R\int(jg_R+\varepsilon-S_j\tau_R)_+\dd\nu.
\end{equation}
For the source in Theorem~\ref{thm:marked}, $G$ is independent of $r$,
so the same formula holds with the correlated weight
$\exp(q\sum_{i=0}^j f_R(B_R^iy))$ inside the integral. For tests
involving collision times or physical samples, the inner $r$ integral
in \eqref{eq:record-flux} must be retained.

\section{Localization to actual alternating itineraries}
\label{sec:geometry}
The smallest distance between distinct disks is $g_R$, attained only
by the six normal chords between nearest neighbors. The next center
distance is $\sqrt3$. An outgoing ray cannot return to the same convex
disk without another reflection.

\begin{lemma}[Uniform localization]\label{lem:localization}
There are a transverse coordinate radius $a_0>0$ and an excess
$\varepsilon_0>0$, independent of $j\ge1$ and $R\in K$, with the
following property. Every regular $j$-flight segment satisfying
$S_j\tau_R<jg_R+\varepsilon_0$ alternates between one nearest-neighbor
pair. Its impacts lie in the facing arcs $|y_i|<a_0$ about the normal
chord. Every segment in these arcs which satisfies the local specular
stationarity equations is an actual, unobstructed billiard segment.
\end{lemma}
\begin{proof}
Every flight has length at least $g_R$. If their sum has excess less
than $\varepsilon_0$, each individual length is at most
$g_R+\varepsilon_0$. Its target center $e$ satisfies
$|e|\le2R+g_R+\varepsilon_0=1+\varepsilon_0$; taking
$\varepsilon_0<\sqrt3-1$ forces $|e|=1$.

For nearest disks, the minimum of the distance between their boundary
points is unique. By compactness in $R\in K$, chords whose length is
within $\varepsilon_0$ of $g_R$ have both endpoints and direction
uniformly close to that normal chord. This compactness statement can
also be read from the endpoint cosine bound
\[
 n_0\cdot v,\ -n_1\cdot v
 \ge\frac{1-2R-\tau^2}{2R\tau},
\]
whose right side tends to one as $\tau\downarrow g_R$, uniformly on
$K$. Reflection at the target therefore sends the direction close to
$-e$. Distinct nearest-neighbor directions are separated by $\pi/3$.
Taking the neighborhoods sufficiently small forces the next short
flight to go back to the preceding disk. This argument uses the same
neighborhood at every step and hence does not deteriorate with $j$.

The closed normal chord between the pair has positive distance from
all other disks, uniformly on $K$. For the pair $(0,(1,0))$, the other
nearest centers have vertical coordinate $\sqrt3/2$, while
$R\le R_+<1/2$; more distant centers are farther away. Small facing
arcs preserve this separation. Their connecting chords are outgoing
and incoming with normal cosines bounded away from zero. Convexity
prevents a second crossing of the originating disk before the target.
Thus these are first physical hits, and the stationarity equations,
which express the equality of tangential velocity components, are
exactly the specular reflection law. There is no grazing cutoff on a
preselected preparation: the localization follows from the complete
event's excess constraint.
\end{proof}

Fix the pair $(0,(1,0))$ and use a global transverse coordinate $y$.
The facing boundary points are
$(\sqrt{R^2-y^2},y)$ on the first disk and
$(1-\sqrt{R^2-y^2},y)$ on the second. The same symmetric length function
applies in both directions:
\begin{equation}\label{eq:ell}
 \ell_R(u,v)=
 \sqrt{\bigl(1-\sqrt{R^2-u^2}-\sqrt{R^2-v^2}\bigr)^2+(v-u)^2}.
\end{equation}
For endpoints $u=y_0$, $v=y_j$, define the local broken-length action
\begin{equation}\label{eq:action}
 \mathcal L_{j,R}(y_0,\ldots,y_j)=\sum_{i=0}^{j-1}\ell_R(y_i,y_{i+1}).
\end{equation}
The interior variables satisfy $\partial_{y_i}\mathcal L_{j,R}=0$.
We write $W_{j,R}(u,v)$ for the stationary value, whose existence,
uniqueness and uniform estimates are proved next.

The length has the even analytic expansion
\begin{equation}\label{eq:one-quadratic}
 \ell_R(u,v)=g_R+
 \frac{c_R(u^2+v^2)-2uv}{2g_R}+O((u^2+v^2)^2).
\end{equation}
All derivatives of the remainder have uniform analytic bounds in a
fixed small neighborhood for $R\in K$. In particular,
$-\partial_{uv}\ell_R(0,0)=1/g_R$. There is also the elementary bound
\begin{equation}\label{eq:energy-coercive}
 \ell_R(u,v)-g_R\ge
 R-\sqrt{R^2-u^2}+R-\sqrt{R^2-v^2}
 \ge\frac{u^2+v^2}{2R}.
\end{equation}
Consequently any localized orbit of total excess $\varepsilon$ satisfies
$y_0^2+y_j^2+2\sum_{i=1}^{j-1}y_i^2\le2R\varepsilon$.
This bounds all interior variables at once, without a factor $j$.

\section{A boundary-value calculus uniform in the itinerary length}
\label{sec:action}
We give the uniform estimates explicitly. In particular, an exponentially
small mixed derivative will be estimated relatively, not as the difference
of two quantities with a fixed absolute error.

\begin{lemma}[Uniform stationary segment]\label{lem:bridge}
For endpoints $|u|,|v|<a_1$, with $a_1>0$ independent of $j$, the
interior equations of \eqref{eq:action} have a unique small solution.
The solution and $W_{j,R}$ are real analytic in $R,u,v$. There are
$C<\infty$ and $0<\rho<1$, independent of $j$, such that, writing
$a=|u|+|v|$,
\begin{align}\label{eq:bridge-estimates}
 y_i&=\frac{\sinh((j-i)\chi_R)}{\sinh(j\chi_R)}u
       +\frac{\sinh(i\chi_R)}{\sinh(j\chi_R)}v
       +O\bigl(a^3(\rho^i+\rho^{j-i})\bigr),\\
 |y_i|&\le Ca(\rho^i+\rho^{j-i}),\qquad 0\le i\le j.\notag
\end{align}
The analytic bounds for the bridge and for $W_{j,R}-jg_R$, and
every fixed parameter derivative of these functions, are uniform in $j$. Moreover
\begin{equation}\label{eq:action-jet}
 W_{j,R}(u,v)=jg_R+\tfrac12(u,v)H_{j,R}(u,v)^t+O(a^4),
\end{equation}
with uniform analytic remainder bounds.
\end{lemma}
\begin{proof}
The linearized interior equations are
\[
 y_{i+1}-2c_Ry_i+y_{i-1}=0,
\]
whose endpoint solution is the displayed linear term. On the $j-1$
interior coordinates their Hessian is
\[
 H^0_{\mathrm{int}}=g_R^{-1}
       \operatorname{tridiag}(-1,2c_R,-1).
\]
For $1\le i\le k\le j-1$, its inverse has entries
\begin{equation}\label{eq:green}
 G_{ik}=g_R\frac{\sinh(i\chi_R)\sinh((j-k)\chi_R)}
                    {\sinh\chi_R\sinh(j\chi_R)},
 \qquad G_{ki}=G_{ik}.
\end{equation}
Substitution in the recurrence verifies the inverse: it solves the
homogeneous equation away from $i=k$, vanishes at both endpoints, and
its discrete derivative has the jump required for a unit vector at $k$.
In particular $|G_{ik}|\le C e^{-\chi_*|i-k|}$.

Choose $\rho\in(e^{-\chi_*},1)$ and put
$w_i=\rho^i+\rho^{j-i}$. A geometric-series calculation gives
\begin{equation}\label{eq:weighted-green}
 \sum_k e^{-\chi_*|i-k|}w_k\le C_\rho w_i.
\end{equation}
For example, for the part $\rho^k$, split the sum into $k\le i$ and
$k>i$ and sum the ratios $e^{-\chi_*}/\rho$ and
$e^{-\chi_*}\rho$, respectively; the other part follows by reversal.
Thus $G$ is uniformly bounded on the weighted supremum norm
$\|z\|_w=\max_i|z_i|/w_i$.

Let $y^{\rm lin}$ be the linear bridge. The nonlinear equations can be
written
$y=y^{\rm lin}-G\mathcal N_R(y)$, where the $i$th coordinate of
$\mathcal N_R$ depends only on $y_{i-1},y_i,y_{i+1}$ and is cubic at zero.
This follows by differentiating the even expansion
\eqref{eq:one-quadratic}. On $\|y\|_w\le C_0a$, the weighted norm of
$\mathcal N_R(y)$ is at most $C a^3$, and its Lipschitz constant in that
ball is at most $C a^2$. Adjacent weights have ratios bounded by
$\rho^{-1}$, so these constants do not depend on the number of variables.
Equations \eqref{eq:green}--\eqref{eq:weighted-green} now give a contraction
for sufficiently small $a_1$, proving existence and \eqref{eq:bridge-estimates}.

The action is strictly convex in the small coordinate box. Its interior
Hessian has diagonal-minus-absolute-off-diagonal row sums at least
$2/R-Ca_0^2$; endpoint rows of the full Hessian have the corresponding
lower bound $1/R-Ca_0^2$. Decreasing the fixed box makes both positive.
Strict convexity proves uniqueness throughout that box, not just in
the contraction ball.

All the estimates also hold on a common small complex neighborhood of
$K$: choose the branches of $\chi_R$ and the square roots there, keep
$\Re\chi_R$ positive, and choose $\rho$ with a strict margin in
\eqref{eq:weighted-green}. The analytic contraction converges uniformly.
Cauchy's estimates on a smaller neighborhood give uniform bounds for
each prescribed derivative. This also justifies differentiating the
bridge with respect to the radius rather than holding an unspecified
collision map fixed.

Finally, $\sum_iw_i^k$ is uniformly bounded for every fixed $k>0$.
The sum of the quartic remainders in \eqref{eq:one-quadratic} is therefore
$O(a^4)$, uniformly in $j$. The linear bridge is stationary for the
quadratic action with its endpoints fixed, so replacing it by the
nonlinear bridge changes that quadratic action by $O(a^6)$.
This proves \eqref{eq:action-jet}, including its analytic bounds.
The case $j=1$ has no interior variables and follows directly from
\eqref{eq:ell}.
\end{proof}

\begin{lemma}[Effective Hessian and twist]\label{lem:hessian}
The endpoint Hessian and mixed derivative at the normal segment are
\begin{equation}\label{eq:hessian}
 H_{j,R}=\frac{s_R}{g_R}
 \begin{pmatrix}\coth(j\chi_R)&-\csch(j\chi_R)\\
                -\csch(j\chi_R)&\coth(j\chi_R)\end{pmatrix},
 \qquad
 d^0_{j,R}:=-W_{j,R,uv}(0,0)=\frac{s_R}{g_R\sinh(j\chi_R)}.
\end{equation}
In particular
\begin{equation}\label{eq:hessian-det}
 \sqrt{\det H_{j,R}}=\frac{s_R}{g_R}=\frac1{R\sqrt{g_R}},
 \qquad
 \frac{d^0_{j,R}}{\sqrt{\det H_{j,R}}}=\frac1{\sinh(j\chi_R)}.
\end{equation}
Both eigenvalues of $H_{j,R}$ are bounded above and below by positive
constants independent of $j$ and $R\in K$.
\end{lemma}
\begin{proof}
For the quadratic action the derivative at its first endpoint is
$(c_Ru-y_1^{\rm lin})/g_R$. Insert the bridge formula and use
$c_R\sinh(j\chi_R)-\sinh((j-1)\chi_R)
=s_R\cosh(j\chi_R)$. Reversal gives the other endpoint derivative,
proving \eqref{eq:hessian}. The identity
$\coth^2x-\csch^2x=1$ proves \eqref{eq:hessian-det}.
The eigenvalues are
$(s_R/g_R)\tanh(j\chi_R/2)$ and
$(s_R/g_R)\coth(j\chi_R/2)$. Their asserted bounds follow from
$j\ge1$, $\chi_R\ge\chi_*>0$ and compactness of $K$.
\end{proof}

An absolute bound $W_{uv}=-d^0_{j,R}+O(a^2)$ would be useless for large
$j$. The next estimate is the relevant replacement.

\begin{lemma}[Relative flux determinant]\label{lem:relative}
On a fixed endpoint neighborhood,
\begin{equation}\label{eq:relative}
 -W_{j,R,uv}(u,v)=d^0_{j,R}\,b_{j,R}(u,v),\qquad
 b_{j,R}=1+O(u^2+v^2)>0.
\end{equation}
The functions $b_{j,R}$ are even under $(u,v)\mapsto(-u,-v)$ and
have uniform analytic bounds, including every fixed radius and endpoint
derivative, independently of $j$.
\end{lemma}
\begin{proof}
At the stationary path put
$b_i=-\ell_{R,uv}(y_i,y_{i+1})>0$ and let $H_{\mathrm{int}}$ be the
interior Hessian of the action. Schur complementation, or the corner
cofactor of this tridiagonal matrix, gives the exact identity
\begin{equation}\label{eq:twist-product}
 -W_{j,R,uv}=\frac{\prod_{i=0}^{j-1}b_i}{\det H_{\mathrm{int}}}.
\end{equation}
For $j=1$ the empty determinant is one. For $j\ge2$, the endpoint-to-interior
entries are $-b_0$ and $-b_{j-1}$. The corner entry of
$H_{\mathrm{int}}^{-1}$ is
$b_1\cdots b_{j-2}/\det H_{\mathrm{int}}$, proving
\eqref{eq:twist-product} with the indicated sign.

The evenness of \eqref{eq:ell} and \eqref{eq:bridge-estimates} imply
\[
 b_i=g_R^{-1}\bigl(1+O(y_i^2+y_{i+1}^2)\bigr),\qquad
 \sum_i|\log(g_Rb_i)|\le C a^2.
\]
Write $H_{\mathrm{int}}=H^0_{\mathrm{int}}+E$. This is a tridiagonal
perturbation. The sum of the absolute values of its entries is at most
$C\sum_i y_i^2\le Ca^2$, and hence its trace norm satisfies
$\|E\|_1\le Ca^2$. Here each matrix unit has trace norm one.
Also $\|(H^0_{\mathrm{int}})^{-1}\|\le C$, either by
\eqref{eq:green} or by the uniform diagonal-dominance margin.
Decreasing the endpoint neighborhood ensures
$\|(H^0_{\mathrm{int}})^{-1}E\|<1/2$. Therefore
\begin{equation}\label{eq:log-det}
 \log\frac{\det H_{\mathrm{int}}}{\det H^0_{\mathrm{int}}}
 =\sum_{m=1}^\infty\frac{(-1)^{m+1}}m
      \tr\bigl[((H^0_{\mathrm{int}})^{-1}E)^m\bigr]=O(a^2).
\end{equation}
Indeed, the absolute trace of the $m$th term is bounded by the
trace norm of its first factor times the operator norm to power $m-1$.
Thus the sum is bounded by $2C\|E\|_1$, not by the dimension of the matrix.

Apply \eqref{eq:twist-product} also at zero and subtract logarithms.
The resulting logarithm of the ratio is analytic and $O(a^2)$ with
constants independent of $j$. The same argument on the complex
neighborhood from Lemma~\ref{lem:bridge}, followed by Cauchy's estimates,
gives all fixed derivative bounds. Exponentiation proves
\eqref{eq:relative}. Simultaneous reversal of all transverse coordinates
leaves the length unchanged and, by uniqueness, reverses the entire
stationary path. It proves the stated evenness. Positivity follows
also directly from \eqref{eq:twist-product} and the positive definite
interior Hessian.
\end{proof}

\begin{proposition}[Physical flux in endpoint coordinates]\label{prop:endpoint-flux}
For each of the six oriented nearest-neighbor segments, the small
endpoint coordinates $(u,v)$ parametrize its outgoing initial collision
state one-to-one, and
\begin{equation}\label{eq:endpoint-flux}
 \dd\nu=\frac{-W_{j,R,uv}(u,v)}{4\pi R}\dd u\dd v.
\end{equation}
Moreover $S_j\tau_R=W_{j,R}(u,v)$ on this chart.
\end{proposition}
\begin{proof}
The first variation of length gives
$\partial_uW=-v_0^+\cdot q_0'(u)$. If $s_0$ is boundary arclength,
the outgoing tangential momentum is
$p_0=-W_u/(\dd s_0/\dd u)$. Thus
$|\dd s_0\dd p_0|=|W_{uv}|\dd u\dd v$.
The mixed derivative is nonzero by Lemma~\ref{lem:relative}, so this
is a local coordinate change. At fixed $u$, $p_0$ is strictly monotone
in $v$ throughout the small box, proving injectivity. The stationary
path is physical by Lemma~\ref{lem:localization}; uniqueness identifies
its successive states with the iterates of $B_R$. Normalization of
$\dd s_0\dd p_0$ by its total mass $4\pi R$ proves the measure formula.
The length identity follows by adding the actual flight lengths.
\end{proof}

\section{Uniform threshold integration and its response}
\label{sec:threshold-proof}
We now integrate the actual flux on the whole onset event. The dimension
of this integration is two, regardless of the number of flights.

\begin{lemma}[Uniform even Morse coordinates]\label{lem:morse}
There is an odd analytic change of variables $z=\Phi_{j,R}(w)$ on a
fixed neighborhood of zero in $\R^2$, with uniform inverse and derivative
bounds, such that
\begin{equation}\label{eq:morse}
 W_{j,R}(\Phi_{j,R}(w))-jg_R=\tfrac12|w|^2,
 \qquad D\Phi_{j,R}(0)=H_{j,R}^{-1/2}.
\end{equation}
In these variables the endpoint flux is
\begin{equation}\label{eq:morse-flux}
 \dd\nu=\frac{1}{4\pi R\sinh(j\chi_R)}
           a_{j,R}(w)\dd w,
 \qquad a_{j,R}(w)=1+O(|w|^2),
\end{equation}
where $a_{j,R}$ is even, positive and uniformly analytic, including
all fixed parameter derivatives.
\end{lemma}
\begin{proof}
Put $E(z)=W_{j,R}(z)-jg_R$. Taylor's integral identity gives
\[
 E(z)=\tfrac12 z^t M(z)z,\qquad
 M(z)=2\int_0^1(1-t)D^2E(tz)\dd t.
\]
The matrix $M$ is symmetric, even and analytic, with $M(0)=H_{j,R}$
and $M(z)-H_{j,R}=O(|z|^2)$ uniformly. The positive definite square
root is analytic on a fixed neighborhood of the uniformly positive
matrices $H_{j,R}$. Thus $w=M(z)^{1/2}z$ is an odd analytic map with
uniformly invertible differential at zero. Its nonlinear differential
is a uniformly small perturbation on a fixed ball. Contraction applied
to the equation $z=H_{j,R}^{-1/2}w+O(|z|^3)$ gives an analytic inverse
on a common smaller ball, with uniform derivative estimates. This proves
\eqref{eq:morse} without a dimension-dependent Morse neighborhood.

Insert \eqref{eq:relative} into \eqref{eq:endpoint-flux} and use
\eqref{eq:hessian-det}. The factor
$b_{j,R}(\Phi(w))\det D\Phi(w)\sqrt{\det H_{j,R}}$ is positive,
even, uniformly analytic and equals $1+O(|w|^2)$. It is the required
$a_{j,R}$. Orientation may be chosen so the local Jacobian is positive.
\end{proof}

The energy bound \eqref{eq:energy-coercive} shows that, for a common
smaller $\varepsilon_0$, the complete sublevel set in each of the six
charts lies inside these coordinates. Conversely all of the small
ball $|w|^2<2\varepsilon$ gives physical segments. Thus no extraneous
cutoff or boundary of a chart enters the following integral.

\begin{proof}[Proof of Theorem~\ref{thm:unweighted}]
By \eqref{eq:Aj} and \eqref{eq:morse-flux}, the contribution of one
orientation is
\[
 \frac{\lambda_R}{4\pi R\sinh(j\chi_R)}
 \int_{|w|^2<2\varepsilon}
      (\varepsilon-|w|^2/2)a_{j,R}(w)\dd w.
\]
The unweighted quadratic integral equals $\pi\varepsilon^2$.
Indeed, rescale $w=\sqrt{2\varepsilon}\,z$ and use
$\int_{|z|<1}(1-|z|^2)\dd z=\pi/2$.
All six orientations have the same flux by lattice symmetry.
The analytic even amplitude has a convergent expansion in homogeneous
even polynomials, with uniform bounds. Integration on the unit disk
therefore gives $1+\varepsilon h_j(R,\varepsilon)$ after division by
the leading term; its analytic bounds are uniform in $j$.
This proves \eqref{eq:main-onset} and every stated derivative bound.
The roof lower bound proves vanishing on the other side of the onset.
\end{proof}

\begin{lemma}[Additive marks on the endpoint bridge]\label{lem:marks}
In the $e$ chart define
\[
 V_{j,e}(R,u,v)=\sum_{i=0}^j
       \bigl(f_R(X_i(R,u,v))-f_R(x_{(-1)^i e})\bigr).
\]
Then $V_{j,e}(R,0,0)=0$, and $V_{j,e}$ and every fixed derivative in
$R,u,v$ have bounds independent of $j$. In particular
$|V_{j,e}|\le C(|u|+|v|)$.
\end{lemma}
\begin{proof}
Each impact position depends on $y_i$. The direction of its adjacent
flight is the chord divided by \eqref{eq:ell}; its denominator is at
least $\min_Kg_R$. Interior and final outgoing velocities are obtained
by the specular formula $v^+=v^--2(v^-\cdot n)n$. These maps are
analytic on the fixed arcs. The last velocity uses $(y_{j-1},y_j)$;
no additional endpoint is needed.

Consequently the deviation of $X_i$ from its normal state is bounded
by $Ca(w_{i-1}+w_i+w_{i+1})$, omitting unavailable indices. The same
weighted bounds hold for every fixed derivative of these deviations:
use the complex estimates in Lemma~\ref{lem:bridge}, or differentiate
the local maps and the bridge. The normal states in angular coordinates
are independent of $R$. Applying the finite chain rule to $f_R$ and
subtracting its value at the normal state gives these weighted bounds
for the summands as well. Since $\sum_iw_i\le2/(1-\rho)$, their sum
has uniform derivative bounds. The value at zero and the linear bound
then follow from the mean value formula.
\end{proof}

\begin{proof}[Proof of Theorem~\ref{thm:marked}]
In one orientation, factor the source into
$e^{qw_{j,e}(R)}e^{qV_{j,e}(R,u,v)}$. Lemmas~\ref{lem:morse} and
\ref{lem:marks} reduce its normalized threshold integral to
\begin{equation}\label{eq:normalized-source}
 F_{j,e}(R,q,\varepsilon)=\frac2\pi
 \int_{|z|<1}(1-|z|^2)
 a_{j,R}(\sqrt{2\varepsilon}\,z)
 \exp\bigl(qV_{j,e}(R,\Phi_{j,R}(\sqrt{2\varepsilon}\,z))\bigr)\dd z.
\end{equation}
It equals one at $\varepsilon=0$. Replace the smooth amplitude in this
integral by its even part in $z$. The integral is unchanged, and every
odd Taylor polynomial disappears. Taylor's formula to any prescribed
finite order, including its differentiated remainder, shows that
$F_{j,e}=1+\varepsilon h_{j,e}$ with all fixed mixed derivative bounds
uniform in $j$. More explicitly, to control $c$ derivatives in
$\varepsilon$ use the even Taylor expansion to degree $2c+4$ on a
slightly larger ball; its integral remainder after the substitution
$w=\sqrt{2\varepsilon}z$ is $O(\varepsilon^{c+2})$ with the
corresponding differentiated bounds. Radius and source derivatives
obey the same argument by Lemma~\ref{lem:marks}. This is finite-order
smooth calculus and makes no analyticity assumption on $f_R$.
Multiplying back the constants in \eqref{eq:morse-flux} proves
\eqref{eq:marked-onset}.

For every fixed $a$, direct differentiation of
$\csch(j\chi_R)=2e^{-j\chi_R}/(1-e^{-2j\chi_R})$ gives
\[
 |\partial_R^a\csch(j\chi_R)|
 \le C_a(1+j)^a e^{-j\chi_*}.
\]
Also $|\partial_R^a w_{j,e}|\le C_a(j+1)$ and
$|e^{qw_{j,e}}|\le e^{q_0F_*(j+1)}$.
The finite product and chain rules, together with the uniform bounds
for $h_{j,e}$, prove \eqref{eq:summable}. Polynomially weighted
geometric series converge, proving termwise differentiation and the
source-vector version.
\end{proof}

\subsection{The singularity across the onset surface}
Set
\[
 C_j(R,q)=\frac{\lambda_R}{4R\sinh(j\chi_R)}
                    \sum_{|e|=1}e^{qw_{j,e}(R)}.
\]
This is positive for real sources. Formula \eqref{eq:marked-onset},
with zero on the negative side, is a full expansion of the form
$\varepsilon_+^2[C_j(R,q)+O(\varepsilon)]$ with smooth right-hand
coefficients. Thus it is $C^1$, and its second $t$ derivative has jump
$2C_j(R,q)$. It is not $C^2$ across any of these onsets. At fixed
physical $t$ and a radius $R_c$ with $t=jg_{R_c}$, the jump of the
second radius derivative is
\begin{equation}\label{eq:radius-jump}
                 8j^2 C_j(R_c,q).
\end{equation}
Indeed $\varepsilon=t-jg_R=2j(R-R_c)$; all other second-derivative
terms have a factor $\varepsilon$ and vanish at the surface.
For higher distributional derivatives, differentiation of
$\varepsilon_+^2$ supplies the corresponding delta terms; for example
its third $t$ derivative contributes $2C_j(R,q)\delta(t-jg_R)$.
At fixed $t$, radius differentiation is
$\partial_R+2j\partial_\varepsilon$ in onset coordinates. Therefore
all fixed-order coefficient and one-sided derivative estimates remain
polynomially weighted, exponentially summable in $j$ in the source
range of \eqref{eq:summable}. This statement concerns the onset
collars, not singularities of other itinerary families at other times.

\section{The joint impact record and physical sampling cuts}
\label{sec:records}
The endpoint parametrization gives more than a count or a marginal roof.
It reconstructs all positions, directions and collision times in the
window. We describe its onset law explicitly.

Let
\begin{equation}\label{eq:paraboloid}
 \mathcal U=\{(z,\eta)\in\R^2\times\R:
                   |z|<1,\ 0<\eta<1-|z|^2\},\qquad
 \dd\mathsf m=\frac2\pi\one_{\mathcal U}\dd z\dd\eta.
\end{equation}
This is a probability measure. For $z\in\R^2$, set
$(U,V)^t=H_{j,R}^{-1/2}z$ and define
\begin{align}\label{eq:limiting-bridge}
 Y_i(z)&=\frac{\sinh((j-i)\chi_R)}{\sinh(j\chi_R)}U
       +\frac{\sinh(i\chi_R)}{\sinh(j\chi_R)}V,\qquad 0\le i\le j,\\
 Q_i(z)&=\frac{Y_i^2+Y_{i+1}^2}{R}
                    +\frac{(Y_{i+1}-Y_i)^2}{g_R},\qquad 0\le i<j.\notag
\end{align}
These quadratic functions satisfy $Q_i\ge0$ and
$\sum_{i=0}^{j-1}Q_i=|z|^2$.

\begin{theorem}[Marked onset record]\label{thm:record}
Condition on $N_{jg_R+\varepsilon}=j+1$. The orientation $e$ is exactly
uniform on the six nearest lattice vectors. In that orientation, use
the exact Morse coordinates $w$ of the endpoints and the first-impact
time $r$, and set $z=w/\sqrt{2\varepsilon}$, $\eta=r/\varepsilon$.
Their law has total variation distance at most $C\varepsilon$ from
$\mathsf m$, uniformly in $j\ge1$ and $R\in K$.

On these same coordinates, the transverse impact positions and impact
times $t_i$ satisfy
\begin{equation}\label{eq:record-asymptotic}
 \frac{y_i}{\sqrt{2\varepsilon}}=Y_i(z)+O(\varepsilon),\qquad
 \frac{t_i-ig_R}{\varepsilon}
       =\eta+\sum_{h=0}^{i-1}Q_h(z)+O(\varepsilon),
 \qquad 0\le i\le j,
\end{equation}
with errors uniform in the index, itinerary length, radius and
$(z,\eta)\in\mathcal U$. Both velocity traces are obtained from the
same bridge by the chord and reflection formulas.

Consider any fixed finite number of samples at
$t^{\rm s}=kg_R+\kappa\varepsilon\in[0,jg_R+\varepsilon]$, with
$0\le k\le j$ and bounded fixed $\kappa$. In the limiting record, the
sample precedes impact $k$ exactly on
\begin{equation}\label{eq:sample-cut}
 \eta+\sum_{h<k}Q_h(z)>\kappa.
\end{equation}
The joint probabilities of the resulting before/after cuts converge
with error $O(\varepsilon)$, uniformly in $j$ and $R$. Thus the
sampling-time interfaces are retained, not suppressed by pointwise
differentiation away from event times.
\end{theorem}
\begin{proof}
Keep the $r$ integral in \eqref{eq:record-flux}. The changes
$w=\sqrt{2\varepsilon}z$, $r=\varepsilon\eta$ give exactly the domain
$\mathcal U$ and, up to its normalizing constant, density
$a_{j,R}(\sqrt{2\varepsilon}z)$. By \eqref{eq:morse-flux}, this density
is $1+O(\varepsilon)$ uniformly. Since
$\operatorname{vol}\mathcal U=\pi/2$, normalization proves the total
variation assertion. Rotation through multiples of $\pi/3$ preserves
phase volume and the event, proving the exact orientation weights.

The odd Morse inverse has expansion
$\Phi_{j,R}(w)=H_{j,R}^{-1/2}w+O(|w|^3)$. Insert it into
\eqref{eq:bridge-estimates}. This proves the first formula in
\eqref{eq:record-asymptotic}, with the stronger error
$O(\varepsilon(\rho^i+\rho^{j-i}))$. Expanding each chord with
\eqref{eq:one-quadratic} gives
$\ell_R(y_i,y_{i+1})=g_R+\varepsilon Q_i(z)
 +O(\varepsilon^2(\rho^i+\rho^{j-i})^2)$.
The sum of these error bounds is uniform, proving the time formula.
The identity for $\sum Q_i$ is the quadratic endpoint action evaluated
at $H_{j,R}^{-1/2}z$.

For clarity, in the pair's fixed global transverse direction the leading
outgoing transverse velocity at impacts $i<j$, divided by
$\sqrt{2\varepsilon}$, is $(Y_{i+1}-Y_i)/g_R$.
At the final impact it is
$(Y_j-Y_{j-1})/g_R+2Y_j/R$; the initial incoming value is
$(Y_1-Y_0)/g_R-2Y_0/R$.
The longitudinal components and the unscaled velocities are given
exactly by the unit chord and specular formulas. Uniform analyticity
on the fixed arcs gives their corresponding Taylor bounds. Together
with the initial and terminal pieces in Proposition~\ref{prop:marked-flux},
this reconstructs the entire physical path.

Consecutive impacts are separated by at least $g_R$, so a sample within
$O(\varepsilon)$ of $kg_R$ can cross only the $k$th impact for small
$\varepsilon$. Formula \eqref{eq:record-asymptotic} gives
\eqref{eq:sample-cut}. The cut has derivative one in $\eta$.
For each fixed $z$, the set where its left side lies within $\delta$
of $\kappa$ has $\eta$-length at most $2\delta$. Its
$\mathsf m$-mass is therefore at most $4\delta$. Taking
$\delta=C\varepsilon$ and adding the $O(\varepsilon)$ density error
bounds the mismatch probability. A union bound proves the result for
any fixed finite set of cuts. Bounded Lipschitz tests of the scaled
positions and times can be added using their uniform pathwise errors.
No differentiability of a discontinuous sampling decoder is assumed.
\end{proof}

For the source-weighted conditional law, the orientation weights tend
to $e^{qw_{j,e}}/\sum_{e'}e^{qw_{j,e'}}$. Given an orientation, the
limiting law is still \eqref{eq:paraboloid}; its total variation error
is $O(\sqrt\varepsilon)$ uniformly for bounded sources, by
Lemma~\ref{lem:marks}. The integrated source normalization has the
stronger $O(\varepsilon)$ relative error of Theorem~\ref{thm:marked}
because its odd part cancels. The distinction between these two errors
is important for a full marked law.

\section{Mechanical consequences of the threshold coefficients}
\label{sec:consequences}
\subsection{A multi-roof effect not determined by the one-roof law}
\begin{corollary}[The nonrenewal onset exponent]\label{cor:nonrenewal}
Fix $R\in(0,1/2)$. Let an independent-roof stationary suspension have
the same section roof marginal as the Lorentz gas. Its intensity is
therefore the same $\lambda_R$. For each fixed $j\ge1$, its maximal
count onset satisfies
\begin{equation}\label{eq:renewal}
 \Pr^{\rm ind}(N_{jg_R+\varepsilon}=j+1)
 =\frac{\lambda_R}{(j+1)!}
       \left(\frac3{\sqrt{g_R}}\right)^j
       \varepsilon^{j+1}\bigl(1+O_j(\varepsilon)\bigr).
\end{equation}
For the actual Lorentz gas the exponent is two for every $j$, with
coefficient \eqref{eq:main-onset}. In particular the two laws have
different leading powers for every $j\ge2$, despite their identical
one-roof marginal and intensity.
\end{corollary}
\begin{proof}
The case $j=1$ of the action calculation gives
$\int(g_R+\varepsilon-\tau_R)_+\dd\nu
=3\varepsilon^2/(2\sqrt{g_R})+O(\varepsilon^3)$.
The smooth threshold density proved above permits differentiation.
Thus the excess roof $u=\tau_R-g_R$ has density
$3/\sqrt{g_R}+O(u)$ near $u=0$.
For $j$ independent excess roofs, integrate their product density on
$u_i\ge0$, $\sum_i u_i<\varepsilon$ with weight
$\varepsilon-\sum_i u_i$. The leading integral is
$\varepsilon^{j+1}/(j+1)!$, as follows by integrating successively
on this simplex. The error is $O_j(\varepsilon)$ relative to it.
The stationary identity \eqref{eq:tails-new}, with
$\varepsilon<g_R$, supplies the factor $\lambda_R$ and excludes the
next count. This proves \eqref{eq:renewal}. The actual law is
Theorem~\ref{thm:unweighted}.
\end{proof}

For example the first genuinely multi-roof count has
\begin{equation}\label{eq:three-hit}
 \Pr_R(N_{2g_R+\varepsilon}=3)
 =\frac{3\lambda_R R}{4(1-R)\sqrt{g_R}}
       \varepsilon^2\bigl(1+O(\varepsilon)\bigr).
\end{equation}
The independent-roof substitute instead has coefficient
$3\lambda_R/(2g_R)$ multiplying $\varepsilon^3$.
The stationary identity is the same in both systems; the difference
comes from the specular boundary-value geometry of $S_2\tau_R$.

\subsection{Recovery of periodic-orbit instability}
\begin{corollary}[Instability in count probabilities]\label{cor:multiplier}
The unstable multiplier of the normal period-two return map is
$e^{2\chi_R}$. The measured threshold coefficients satisfy
\begin{equation}\label{eq:ratio}
 \lim_{j\to\infty}\frac{C_{j+1}(R)}{C_j(R)}=e^{-\chi_R},\qquad
 \frac{C_1(R)}{C_2(R)}=2\cosh\chi_R.
\end{equation}
Thus they determine that multiplier. The unweighted coefficient series
has the exact representation
\begin{equation}\label{eq:lambert}
 \sum_{j\ge1}C_j(R)z^j
 =\frac{3\lambda_R}{R}
       \sum_{m\ge0}\frac{z e^{-(2m+1)\chi_R}}
                         {1-z e^{-(2m+1)\chi_R}},\qquad |z|<e^{\chi_R}.
\end{equation}
It extends meromorphically to the complex plane with simple poles
at $z=e^{(2m+1)\chi_R}$, $m\ge0$.
\end{corollary}
\begin{proof}
The Jacobi recurrence in the proof of Lemma~\ref{lem:bridge} has
transfer matrix $\begin{psmallmatrix}2c_R&-1\\1&0\end{psmallmatrix}$
and eigenvalues $e^{\chi_R},e^{-\chi_R}$. The local change between
successive transverse positions and a collision position--momentum
pair is invertible since the one-flight twist is $1/g_R$.
After two flights the coordinate frame returns to the original
orientation, so the physical period-two derivative is conjugate to
the square of this transfer matrix. This proves the multiplier claim.
The ratios follow from the explicit $C_j$.

Expand $\csch(j\chi)=2\sum_{m\ge0}e^{-(2m+1)j\chi}$ and sum the
absolutely convergent geometric series in $j$. On every compact set
avoiding the displayed poles, the tail in $m$ converges uniformly,
since its terms are $O(|z|e^{-(2m+1)\chi})$. This gives the
meromorphic continuation. At each indicated point only one summand
has a pole and its residue is nonzero, so there is no cancellation.
\end{proof}

Series \eqref{eq:lambert} organizes threshold coefficients belonging to
different observation windows. It is not the ordinary pressure or the
generating function of counts in one fixed window. The recovery of
instability here uses equilibrium tail amplitudes; it does not claim
that the relation between periodic action Hessians and monodromy is new.

\subsection{The retained regular-window response}
The complete companion \cite{TC} proves its fixed-window theorem,
including the all-order moving-level formula, without using the present
threshold asymptotic. We record the terminal-level witness requested
for that calculation. At $T=11/100$, $R\in[9/20,47/100]$, choose
$\alpha=0$ facing $(1,0)$ and
\begin{equation}\label{eq:terminal-witness}
 \cos\vartheta=\frac{1-2R+T^2}{2T(1-R)}.
\end{equation}
The right side lies strictly between zero and one. The point
$(R,0)+T(\cos\vartheta,\sin\vartheta)$ lies on the target circle,
since substitution gives squared distance $R^2$ from its center.
It is an incoming hit because
\[
 (1-R)\cos\vartheta-T=\frac{1-2R-T^2}{2T}>0.
\]
The outgoing sign and the no-obstruction estimate in the companion
make it a genuine first hit. The terminal level is therefore nonempty.
Its uniform gradient bound and $\partial_R\tau<0$ show directly that
the level integral in the companion's second-derivative formula is
strictly positive. This is compatible with the onset singularity:
the fixed window remains strictly above $g_R$ on its specified interval,
whereas \eqref{eq:main-onset} crosses the critical level itself.

\section{An open geometric class and complete-event localization}
\label{sec:g-localization}
We verify the geometric assertions used above; in particular, minimizer
selection is not an additional untested hypothesis. Boundary regularity
and all parameter estimates in this section refer to the fixed family
specified in Section~\ref{sec:geometric-results}.

\begin{lemma}[Persistence of the six closest channels]\label{lem:g-channels}
For a sufficiently small $C^4$ neighborhood of a circle, the translates
of $D_\xi$ are disjoint, have uniformly positive curvature, and leave
positive free area. The only pairs capable of a flight within a fixed
positive excess of the shortest gap are those with displacement in
$\mathcal E$. Each such pair has a unique closest chord, with smooth
endpoints, positive nondegenerate distance Hessian and a uniform
clearance from all other scatterers.

There exist fixed contact arcs and $\varepsilon_0>0$ such that every
$j$-flight segment with total length less than $jg_*+\varepsilon_0$
alternates in one of these channels. Its impacts lie in those arcs.
The same constants work for all $j\ge1$.
\end{lemma}
\begin{proof}
A support function gives the boundary parametrization
$q(\theta)=h(\theta)n(\theta)+h'(\theta)t(\theta)$, whose derivative
is $(h+h'')t$. Small $C^2$ distance from $R_*$ keeps $h+h''$ and the
separation $1-2\sup h$ positive. The remaining lattice vectors have
length at least $\sqrt3$; their boundary distances are at least
$\sqrt3-2\sup h$, whereas the candidate gaps are close to $1-2R_*$.
This gives a strict uniform gap to all other translates. Area and
clearance persist by compactness and small Hausdorff distance.

At a closest pair, choose a common transverse coordinate, translating
and rotating its chord to the horizontal segment $[0,g]$. The facing
arcs are $(-\psi_0(u),u)$ and $(g+\psi_1(v),v)$, where
$\psi_b(0)=\psi_b'(0)=0$ and $\psi_b''(0)=\kappa_b>0$. Their distance
has Hessian
\begin{equation}\label{eq:g-single-hessian}
 g^{-1}\begin{pmatrix}1+g\kappa_0&-1\\-1&1+g\kappa_1\end{pmatrix}.
\end{equation}
Its determinant is positive and its eigenvalues have uniform positive
bounds. The implicit function theorem gives smooth continuation of the
circular closest pair. The minimum remains there globally: away from
fixed neighborhoods of the circular endpoints the distance exceeds its
minimum by a fixed amount. Alternatively, the shortest displacement
between two convex bodies is unique, and strict convexity makes each
support point unique. The local Hessian proves the required uniform
nondegeneracy. A finite choice of oriented frames near the six circular
normals makes all graph coordinates smooth in $\xi$.

Every flight has length at least $g_*$. If the total excess is less
than $\varepsilon_0$, every individual flight has length at most
$g_*+\varepsilon_0$. Its pair belongs to $\mathcal E$ by the lattice
gap. Its length exceeds that pair's own minimum by at most
$\varepsilon_0$, so compactness and the distance Hessian put its
endpoints and direction in the chosen normal neighborhoods. At the
terminal point, specular reflection reverses a nearly normal direction.
Different candidate directions are still separated by a fixed positive
angle. The next short flight must therefore return to the preceding
translate, and cannot switch to another candidate. This step has no
length-dependent constant. Convexity and the clearance of the normal
chord ensure that each local connecting chord is a first hit, with
normal cosines uniformly away from zero. Tangencies and their regular
preimages form null sets and do not change the full-preparation event.
\end{proof}

For one channel suppress its subscript and write $g=g_e$, $\kappa_b=
\kappa_{e,b}$. The exact successive lengths are
\begin{equation}\label{eq:g-length}
 \ell_b(u,v)=\sqrt{(g+\psi_b(u)+\psi_{1-b}(v))^2+(v-u)^2},
 \qquad b=0,1,
\end{equation}
where the index alternates and the same global transverse direction is
used at both ends. The expansion and coercivity are
\begin{align}\label{eq:g-length-jet}
 \ell_b(u,v)&=g+\frac{c_bu^2+c_{1-b}v^2-2uv}{2g}
                      +O((|u|+|v|)^3),\\
 \ell_b(u,v)-g&\ge\psi_b(u)+\psi_{1-b}(v)
                      \ge c_{\mathrm{loc}}(u^2+v^2).\label{eq:g-coercive}
\end{align}
Here $c_{\mathrm{loc}}>0$ is a uniform coercivity constant. On a fixed small box the full broken-length
Hessian is strictly positive: its endpoint diagonal-dominance margins
are $\kappa_b$ at zero and its interior margins $2\kappa_b$.
Perturbation of each row is bounded by the maximum neighboring coordinate,
not by the number of rows.

\section{Unequal curvatures and relative flux at arbitrary length}
\label{sec:g-action}
For endpoints $y_0=u,y_j=v$, let $W_{j,e}(\xi,u,v)$ be the stationary
value of $\sum_{i=0}^{j-1}\ell_{i\bmod2}(y_i,y_{i+1})$.
We prove its existence and estimates rather than assuming a long
initial-value iteration stays in a common chart.

\begin{lemma}[Alternating Jacobi reduction]\label{lem:g-jacobi}
Put $c=\sqrt{c_0c_1}$, $\gamma=\arcosh c$, and
$\sigma_i=\sqrt{c_{1-(i\bmod2)}}$. Let
$D_j=\diag(\sigma_0,\sigma_j)$. The quadratic endpoint Hessian is
\begin{equation}\label{eq:g-effective}
 H_{j,e}=\frac{c\sinh\gamma}{g}\,
 D_j^{-1}
 \begin{pmatrix}\coth(j\gamma)&-\csch(j\gamma)\\
                 -\csch(j\gamma)&\coth(j\gamma)\end{pmatrix}D_j^{-1}.
\end{equation}
With $d^0_{j,e}=-W_{uv}(\xi,0,0)$ this gives
\begin{equation}\label{eq:g-ratio}
 d^0_{j,e}=\frac{c\sinh\gamma}{g\sigma_0\sigma_j\sinh(j\gamma)},
 \qquad \sqrt{\det H_{j,e}}=\frac{c\sinh\gamma}{g\sigma_0\sigma_j},
 \qquad \frac{d^0_{j,e}}{\sqrt{\det H_{j,e}}}=\csch(j\gamma).
\end{equation}
All endpoint eigenvalues are bounded above and below uniformly in $j$.
The period-two return multiplier is $e^{2\gamma}$.
\end{lemma}
\begin{proof}
The quadratic interior equation is
$y_{i+1}-2c_{i\bmod2}y_i+y_{i-1}=0$.
Set $y_i=\sigma_i z_i$. Since $c_{i\bmod2}\sigma_i/
\sigma_{i+1}=c$, the equation becomes
$z_{i+1}-2cz_i+z_{i-1}=0$. The whole quadratic action becomes
\[
 \frac c{2g}\sum_{i=0}^{j-1}
           \{c(z_i^2+z_{i+1}^2)-2z_i z_{i+1}\}.
\]
Solving its Dirichlet recurrence and differentiating its endpoint value
as in Lemma~\ref{lem:hessian} gives \eqref{eq:g-effective}.
The identity $\coth^2-\csch^2=1$ gives \eqref{eq:g-ratio}.
The middle matrix has eigenvalues $\tanh(j\gamma/2)$ and
$\coth(j\gamma/2)$; $\gamma$ has a positive lower bound and
$\sigma_i$ have fixed positive bounds. This proves uniform coercivity.
The change of scale is two-periodic. After two steps the recurrence
transfer matrix is conjugate to the square of
$\begin{psmallmatrix}2c&-1\\1&0\end{psmallmatrix}$.
The one-flight twist is $1/g$, so position-pair coordinates and collision
position--momentum coordinates are locally equivalent. This proves the
physical multiplier assertion.
\end{proof}

\begin{lemma}[Uniform nonlinear bridge and relative determinant]
\label{lem:g-relative}
On one fixed endpoint neighborhood the stationary segment exists uniquely.
It is smooth in $\xi,u,v$. For some $\rho<1$ and $a=|u|+|v|$,
\begin{align}\label{eq:g-bridge}
 y_i^{\rm lin}&=\sigma_i\left(
 \frac{\sinh((j-i)\gamma)}{\sinh(j\gamma)}\frac u{\sigma_0}
 +\frac{\sinh(i\gamma)}{\sinh(j\gamma)}\frac v{\sigma_j}\right),\\
 |y_i|&\le Ca(\rho^i+\rho^{j-i}),\qquad
 |y_i-y_i^{\rm lin}|\le Ca^2(\rho^i+\rho^{j-i}),\notag\\
 W_{j,e}-jg_e&=\tfrac12(u,v)H_{j,e}(u,v)^t+O(a^3),\notag\\
 -W_{uv}(\xi,u,v)&=d^0_{j,e}(\xi)\,b_{j,e}(\xi,u,v),\qquad
 b_{j,e}(\xi,0,0)=1,\quad b_{j,e}=1+O(a)>0.\label{eq:g-relative}
\end{align}
Every fixed derivative of the action remainder and of $b_{j,e}$ is
uniformly bounded independently of $j$. Every fixed derivative of the
bridge coordinates obeys endpoint-decaying bounds (with the powers of
$a$ reduced by endpoint differentiation). These assertions do not require
that either boundary graph be even.
\end{lemma}
\begin{proof}
Let $\gamma_* =\min_{\xi,e}\gamma_e>0$. Scaling the constant-coefficient
Dirichlet inverse gives, for $i\le k$,
\[
 G_{ik}=\frac{g\sigma_i\sigma_k}{c}
    \frac{\sinh(i\gamma)\sinh((j-k)\gamma)}
         {\sinh\gamma\sinh(j\gamma)},\qquad G_{ki}=G_{ik}.
\]
This is checked by its recurrence, zero endpoint values, and unit jump
at $i=k$. It is bounded by $C e^{-\gamma_*|i-k|}$. Choose
$\rho\in(e^{-\gamma_*},1)$. The geometric-series estimate
\eqref{eq:weighted-green} is unchanged. The nonlinear interior gradient
is now quadratic at zero; on the ball with weighted norm at most $Ca$
its norm is $O(a^2)$ and its Lipschitz constant $O(a)$. Adjacent weights
have bounded ratios. The equation $y=y^{\rm lin}-G\mathcal N(y)$ is a
contraction on this ball for a uniform small endpoint radius. The strict
convexity following \eqref{eq:g-coercive} gives uniqueness throughout the
contact box, including any localized physical segment.

We spell out the smooth-parameter point. Derivatives of the Green entries
are finite sums with polynomial factors in the distances of indices.
The endpoint bridge has analogous factors multiplying its endpoint
exponentials. At each prescribed order these are absorbed by choosing a
strictly slower decay rate still below one. Differentiating the fixed
point equation once moves a linear perturbation of norm less than
$1/2$ to the left. Its inverse has norm at most two on the weighted
space. Repeating this argument expresses an order-$m$ derivative in
terms of lower derivatives, the differentiated local nonlinearities,
and this same inverse. Induction proves the stated bounds; only finitely
many higher boundary and parameter norms are used at each order.
Furthermore $\sum_i(\rho^i+\rho^{j-i})^p$ is bounded for every $p>0$.
Summing the cubic remainders in \eqref{eq:g-length-jet} and using
stationarity of the quadratic bridge proves the action estimate and
its fixed-order derivative bounds.

For the relative statement let $b_i=-\ell_{i\bmod2,uv}(y_i,y_{i+1})$
and let $H_{\rm int}=H^0_{\rm int}+E$ be the interior Hessian. The exact
corner-cofactor identity is still
\[
       -W_{uv}=\frac{\prod_{i=0}^{j-1}b_i}{\det H_{\rm int}}.
\]
Here $b_i=g^{-1}(1+O(|y_i|+|y_{i+1}|))$, and
\[
 \sum_i|\log(gb_i)|\le Ca,\qquad
 \|E\|_1\le\sum_{i,k}|E_{ik}|\le Ca.
\]
The second estimate uses tridiagonality and summable endpoint weights.
Also $\|(H^0_{\rm int})^{-1}\|\le C$. Thus the logarithmic determinant
series in \eqref{eq:log-det} converges, and its trace is $O(a)$ without
a dimension factor. Dividing the cofactor identity by its value at zero
and subtracting logarithms proves \eqref{eq:g-relative}.
For derivatives, each differentiated $E$ has bounded trace norm, and
each differentiated inverse of $H^0_{\rm int}$ has bounded operator
norm: use $\partial H^{-1}=-H^{-1}(\partial H)H^{-1}$ and the uniformly
banded derivatives of $H$. Differentiating the convergent trace series,
or its inverse formula, therefore gives the same fixed-order bounds.
The differentiated sum of $\log(gb_i)$ is controlled by the bridge
weights. Positivity follows from positive twists and the positive
interior determinant. For $j=1$ the determinant is empty and the same
statements follow directly from \eqref{eq:g-length}.
\end{proof}

\section{Threshold integration, moving onsets, and source convergence}
\label{sec:g-integration}
Let $P(\xi)$ be the scatterer perimeter. The collision section measure is
$\dd s\dd p/(2P)$, the mean roof $\pi A/P$, and the stationary collision
intensity $P/(\pi A)$. These identities follow from the same flight-tube
calculation as \eqref{eq:santalo}: irrational torus lines hit the open
scatterer, rational directions form a null angular set, and the positive
minimal gap excludes finite-time accumulation. In particular, in
first-impact endpoint coordinates the full phase measure is
\begin{equation}\label{eq:g-full-flux}
             \frac{-W_{uv}}{2\pi A(\xi)}\dd u\dd v\dd r.
\end{equation}
Indeed, first variation gives $p=-W_u/(\dd s/\dd u)$; the arclength
factor cancels in $|\dd s\dd p|=|W_{uv}|\dd u\dd v$. The positive
relative twist makes the endpoint map injective at fixed $u$. The
localized stationary path is physical by Lemma~\ref{lem:g-channels}.
Thus \eqref{eq:g-full-flux} is not a substituted transverse ensemble.

\begin{lemma}[Uniform smooth radial integration]\label{lem:g-radial}
There are smooth Morse inverses $\Phi_{j,e}(\xi,w)$ on a common ball,
with all fixed derivatives bounded uniformly in $j$, such that
\[
 W_{j,e}(\xi,\Phi_{j,e}(\xi,w))-jg_e=|w|^2/2,
 \qquad D_w\Phi_{j,e}(\xi,0)=H_{j,e}^{-1/2}.
\]
The phase density \eqref{eq:g-full-flux} in $(w,r)$ is
\begin{equation}\label{eq:g-morse-density}
 \frac{a_{j,e}(\xi,w)}{2\pi A(\xi)\sinh(j\gamma_e)}\dd w\dd r,
 \qquad a_{j,e}(\xi,0)=1,
 \qquad a_{j,e}=1+O(|w|)>0,
\end{equation}
with the same derivative bounds. If $a$ is any such uniformly smooth
amplitude, then
\[
 \frac2\pi\int_{|z|<1}(1-|z|^2)a(\xi,\sqrt{2d}\,z)\dd z
            =a(\xi,0)+dH(\xi,d),
\]
where $H$ is smooth for $d\ge0$ and has uniform fixed-order bounds.
No evenness of $a$ or of the scatterer is required.
\end{lemma}
\begin{proof}
For $E=W-jg$, Taylor's integral formula writes
$E(y)=\frac12 y^tM(y)y$, where
$M(y)=2\int_0^1(1-s)D^2E(sy)\dd s$ is uniformly positive near zero.
The map $y\mapsto M(y)^{1/2}y$ has uniformly invertible differential at
zero and small nonlinear differential on a common ball. Its inverse is
obtained by contraction, then differentiated to each finite order.
This is the smooth version of Lemma~\ref{lem:morse}. Inserting the
Jacobian and \eqref{eq:g-ratio} into \eqref{eq:g-full-flux} gives
\eqref{eq:g-morse-density}.

For the last assertion replace $a(w)$ by $(a(w)+a(-w))/2$ in the
integral. Every odd Taylor term vanishes. Taylor expansion to any
prescribed order, with a sufficiently high differentiated remainder,
then yields an expansion in integer powers of $d$. One can also set
$r=\sqrt d$ first: the integral is a smooth even function of $r$ with
uniform bounds. Successive division of its even Taylor remainder by
$r^2$ shows that it is a smooth function of $r^2$ on the right.
Using two more endpoint derivatives for each division proves the
uniform bounds for every fixed derivative of $H$.
\end{proof}

\begin{proof}[Proof of Theorem~\ref{thm:g-stability}]
In a channel's window $jg_e+d$, each of its $j$ complete flights is at
least $g_e$. Hence the allowed first residual time satisfies
$0<r<d-(W-jg_e)<d<g_*$. The preceding roof cannot truncate this interval,
and the next roof cannot fit after the last hit. The direct
first-impact formula of Proposition~\ref{prop:marked-flux} therefore
applies also to this selected channel at its own threshold.
The coercivity \eqref{eq:g-coercive} implies
$|y_0|^2+|y_j|^2+2\sum_{0<i<j}|y_i|^2\le Cd$ on its sublevel.
It places the entire active set inside the common Morse coordinates,
so no chart-boundary term occurs. Its exact probability is
\[
 \frac1{2\pi A\sinh(j\gamma_e)}
 \int_{|w|^2<2d}(d-|w|^2/2)a_{j,e}(\xi,w)\dd w.
\]
Since the quadratic integral is $\pi d^2$, Lemma~\ref{lem:g-radial}
proves the channel assertion, including its derivative bounds.

For the complete event let $t=jg_*+\varepsilon$ with
$0<\varepsilon<\varepsilon_0$. No $j+2$ hits are possible, since
$t<(j+1)g_*$. Lemma~\ref{lem:g-channels} partitions every contributing
$j$-flight segment into exactly one of the six oriented channels.
Its available excess there is $d_{j,e}=\varepsilon-j(g_e-g_*)\le
\varepsilon$. Channels with nonpositive excess contribute zero.
Summing the exact channel formulas gives \eqref{eq:g-competition}.
The minimum-roof bound proves the negative-side assertion.
\end{proof}

\begin{proof}[Proof of Theorem~\ref{thm:g-marked}]
Write the mark on a stationary bridge as $w_{j,e}+V_{j,e}$, subtracting
the normal value separately at each impact. Chord directions and normal
reflections are smooth maps of neighboring transverse coordinates, with
uniformly positive chord denominators. Lemma~\ref{lem:g-relative}
therefore gives $|V_{j,e}|\le C(|u|+|v|)$ and uniform fixed-order
bounds for all its derivatives. The normal states now depend on $\xi$,
but each summand's deviation vanishes at $(u,v)=0$ for every $\xi$;
parameter differentiation preserves its endpoint localization.
The finite chain rule and $\sum_i(\rho^i+\rho^{j-i})\le C$ prove
the bounds. Multiplying the Morse amplitude by $e^{qV_{j,e}}$ and
applying Lemma~\ref{lem:g-radial} proves \eqref{eq:g-marked}.

For the convergence assertion, exact parity counting gives
\begin{equation}\label{eq:g-parity}
 w_{j,e}=(j+1)\bar f_e+
       \one_{\{j\text{ even}\}}(a_{e,0}-a_{e,1})/2.
\end{equation}
For fixed real $q$ the parity factor is bounded above and below by
positive constants, as is the normalized integral of $a_{j,e}e^{qV}$
on a sufficiently small collar. Also $\csch(j\gamma_e)\asymp
2e^{-j\gamma_e}$, uniformly in $j$. Thus a channel's normalized term
is comparable to $e^{j(q\bar f_e-\gamma_e)}$. Positivity rules out
cancellation: equality in \eqref{eq:g-source-domain} gives a
nondecaying term, and a positive exponent gives a growing term.
This proves necessity as well as sufficiency, including the right limit
at zero. On a compact subset of the strict domain there is a common
margin $\delta>0$. Parameter and source derivatives of the leading
factor cost only a polynomial in $j$; offset derivatives act on the
uniform smooth remainder. Every fixed derivative is bounded by
$C(1+j)^m e^{-\delta j}$ and is summable. The vector-source assertion
is the same argument with scalar products.
\end{proof}

\begin{corollary}[Moving threshold interfaces]\label{cor:g-interfaces}
Within the collar of \eqref{eq:g-competition}, the probability is a
$C^1$ function of the smooth physical variables $(\xi,t)$, with smooth
pieces separated by $t=jg_e(\xi)$. For an orientation $e$, along a smooth
scalar path $s\mapsto(\xi(s),t(s))$ entering $d_{j,e}>0$ transversely,
its contribution has second-derivative jump from the inactive to the active side
\begin{equation}\label{eq:g-interface-jump}
       \frac{(t'(s)-j\,\partial_s g_e)^2}{A(\xi)\sinh(j\gamma_e)}.
\end{equation}
Coincident contributions, including the reversed orientation, add with
their respective activation signs.
The marked formula has the additional factor $e^{qw_{j,e}}$.
One-sided and distributional higher derivatives follow by differentiating
the positive-part expression, with polynomial-in-$j$ coefficient bounds.
\end{corollary}
\begin{proof}
Each summand is $d_+^2[C(\xi)+d_+C(\xi)H(\xi,d_+)]$.
It and its first derivatives match at zero. Only
$2C(\partial_s d)^2$ survives in the second-derivative jump.
All other terms contain a factor $d$. Finite addition proves the
claim at intersections. For higher derivatives the ordinary product
and chain rules on each side, and the derivatives of $d_+^2$ in the
distributional sense, give the stated formulas; no differentiation of
$g_*$ is used.
\end{proof}

For a concrete change of minimizing mechanism, take
$h_s(\theta)=R_*+s\cos(2\theta)$. Differentiating the closest distance
at the circular normal chord gives
$g_e'(0)=-2\cos(2\arg e)$: the endpoint variation is the normal support
variation, since the unperturbed endpoint gradient is zero.
The horizontal pair has derivative $-2$ and the other two pairs derivative
$1$. Reflection symmetry keeps those two lengths equal. Consequently
the horizontal pair minimizes for small positive $s$, while the other
two minimize for small negative $s$. Formula \eqref{eq:g-competition}
continues through $s=0$. A nonminimal pair contributes exactly while
$j(g_e-g_*)<\varepsilon$; the parameter transition scale is therefore
$\varepsilon/j$ at a transverse length crossing. This assertion is about
activation of physical channels, not a new family of symbolic itineraries.

\section{Curvatures invisible to all threshold locations}
\label{sec:g-inverse}
The following construction separates geometric uncertainty from its
statistical measurement. Set $\theta_r=r\pi/3$, $r=0,1,2$, and let
\begin{equation}\label{eq:g-isogap-family}
 h_{R,\alpha,\beta,\zeta}(\theta)
   =R+(1-\cos6\theta)b(\theta),\qquad
 b(\theta)=\alpha+\beta\cos2\theta+\zeta\sin2\theta.
\end{equation}
Here $R$ ranges over a compact interval in $(0,1/2)$, and the three
other parameters are sufficiently small. These are smooth centrally
symmetric strictly convex scatterers, generally noncircular. At each
$\theta_r$ and its opposite, $h=R$, $h'=0$, and
\begin{equation}\label{eq:g-contact-curvatures}
       \kappa_r^{-1}=h(\theta_r)+h''(\theta_r)
                     =R+36b(\theta_r).
\end{equation}
The map from $(\alpha,\beta,\zeta)$ to the three $b(\theta_r)$ is
invertible: its rows are $(1,1,0)$, $(1,-1/2,\sqrt3/2)$ and
$(1,-1/2,-\sqrt3/2)$.

\begin{theorem}[Equal gaps and identifiable curvature amplitudes]
\label{thm:g-inverse}
In \eqref{eq:g-isogap-family} all six shortest gaps are exactly
$g=1-2R$, independently of $(\alpha,\beta,\zeta)$. The entire set of
maximal-count threshold locations is therefore the same. Nevertheless,
write
\[
 C_j=\lim_{d\downarrow0}d^{-2}\Pr(N_{jg+d}=j+1),\qquad
 \gamma_r=\arcosh(1+g\kappa_r).
\]
Together with the observed spacing $g$, the unlabelled sequence
$(C_j)_{j\ge1}$ determines the unordered triple
$(\kappa_0,\kappa_1,\kappa_2)$, including multiplicities, and the free
area $A$, without prior knowledge of $A$. It obeys
\begin{equation}\label{eq:g-mobius}
 C_j=\frac1A\sum_{r=0}^2\csch(j\gamma_r),\qquad
 F_j:=\sum_{\substack{k\ge1\\ k\text{ odd}}}\mu(k)C_{kj}
       =\frac2A\sum_{r=0}^2e^{-j\gamma_r},
\end{equation}
where $\mu(k)$ is the arithmetic M\"obius function. Both series involved
in the inversion are absolutely convergent.

In the subfamily $\beta=\zeta=0$, two coefficients already determine
its shape parameter:
\begin{equation}\label{eq:g-two-amplitude-inverse}
 \alpha=\frac1{36}\left\{
       \frac{g}{C_1/(2C_2)-1}-R\right\}.
\end{equation}
Threshold locations alone cannot determine this parameter.
\end{theorem}
\begin{proof}
For small parameters $h+h''$ remains positive, which proves smooth
strict convexity and \eqref{eq:g-contact-curvatures}. For each nearest
lattice vector $e$ at angle $\theta_r$, the two supporting lines have
separation $1-h(\theta_r)-h(\theta_r+\pi)=g$.
Every displacement from the first body to its translate has projection
at least $g$ in direction $e$, and hence length at least $g$.
Since $h'=0$ at the two normals, the supporting points themselves lie
on the axis and attain that distance. More distant translates have
distance at least $\sqrt3-2\sup h>g$ after a uniform decrease of the
parameter neighborhood. Thus the global shortest length is exactly
$g$ and all six channels attain it. The three curvatures can vary
independently by the displayed invertible matrix.

Central symmetry makes both curvatures of the $r$th pair equal.
Theorem~\ref{thm:g-stability} gives the first formula in
\eqref{eq:g-mobius}, with two orientations per pair. Put
$x_r=e^{-\gamma_r}\in(0,1)$. Expanding
$\csch(j\gamma_r)=2\sum_{m\ge1,\ m\text{ odd}}x_r^{mj}$
and interchanging the absolutely convergent sums gives
\[
 F_j=\frac2A\sum_r\sum_{\substack{n\ge1\\n\text{ odd}}}
        x_r^{nj}\sum_{k\mid n}\mu(k)=\frac2A\sum_r x_r^j.
\]
Absolute convergence follows, for example, by bounding the double sum
by $\sum_{n\ge1}n x_{\max}^{nj}<\infty$; the number of divisors
is at most $n$. The divisor identity leaves only $n=1$.

For completeness we give a finite-rank recovery from the transformed
sequence, so uniqueness is not merely asserted. Let $x_1,\ldots,x_d$
be its distinct nodes, $d\le3$, with weights $w_l=2m_l/A>0$, where
$m_l$ are their multiplicities. Then $F_j=\sum_l w_lx_l^j$.
The Hankel matrices
\[
 H^{(0)}=(F_{p+q+1})_{0\le p,q<d},\qquad
 H^{(1)}=(F_{p+q+2})_{0\le p,q<d}
\]
satisfy $H^{(0)}=V\diag(w_lx_l)V^t$ and
$H^{(1)}=V\diag(w_lx_l^2)V^t$, where $V_{pl}=x_l^p$.
The rank of the full Hankel sequence is $d$, and $H^{(0)}$ is positive
definite. Thus the generalized eigenvalues of this pair are exactly
$x_l$. A Vandermonde system using $F_1,\ldots,F_d$ recovers the $w_l$.
Since $\sum_lm_l=3$, we recover
$A=6/(\sum_lw_l)$ and then $m_l=Aw_l/2$.
Finally $\kappa_l=(\cosh(-\log x_l)-1)/g$ proves the curvature claim.
No label ordering or global determination of the boundary is asserted.

If $b\equiv\alpha$, the three curvatures coincide and
$C_1/C_2=2\cosh\gamma=2(1+g/(R+36\alpha))$.
Solving gives \eqref{eq:g-two-amplitude-inverse}. All threshold locations
are independent of $\alpha$ on a nontrivial open interval, while this
ratio varies strictly with it. This proves the claimed nonredundancy.
\end{proof}

The infinite original sequence enters the M\"obius sums; this is an exact
identification theorem, not a finite-noisy-data reconstruction guarantee.
The special two-amplitude formula needs no such infinite transform.
In a circle, by contrast, $R=(1-g)/2$ and $1/C_j$ already satisfies a
second-order recurrence. Corollary~\ref{cor:multiplier} remains correct;
Theorem~\ref{thm:g-inverse} supplies the additional geometric uncertainty
which is absent from that one-parameter model.

\begin{corollary}[Instability selection at equal metric thresholds]
\label{cor:g-selection}
In \eqref{eq:g-isogap-family}, conditional on the maximal-count event
in a common small collar, the mass of a channel with
$\gamma_r>\min_s\gamma_s$ is at most
$C\exp[-j(\gamma_r-\min_s\gamma_s)]$, uniformly in the offset.
Thus channels with the same exact threshold length can have exponentially
different probabilities. The limiting leading-amplitude weights among
minimizing channels are equal.
\end{corollary}
\begin{proof}
The relative remainders in \eqref{eq:g-competition} are bounded above
and below by positive constants. Divide a channel's coefficient by
one with minimal $\gamma$ and use the exponential expression for
$\csch$. At offset zero the equal minimal coefficients give equal
leading-amplitude weights. This last assertion concerns amplitudes;
it does not identify the fixed-positive-offset limiting corrections.
\end{proof}

\input{v2/12_record_response}
\section{Comparison and the remaining long-time problem}
\label{sec:comparison}
\subsection{Stationary identities and uniform orbit segments}
The relation between stationary entry times and section return times is
classical. Marklof \cite{Marklof} develops this relation through stationary
point processes and Palm distributions, without an ergodicity assumption.
The suspension identity in Section~\ref{sec:flux} is an elementary version
adapted to a finite physical record. It supplies the preparation factor
and the residual-time integral, not the geometry or mass of its active set.
The new calculation starts with $S_j\tau_R$, rather than with a supplied
independent roof sequence.

The use of length as a billiard generating function, and the relation of
action Hessians to linearized dynamics, likewise belong to classical
variational mechanics. Bolotin and Treschev \cite{Hill} treat Hill-type
determinant identities for discrete and continuous Lagrangian systems,
including symplectic twist maps. Our corner-cofactor identity is a finite
Dirichlet calculation of this kind; it is not claimed as a new general
Hill formula. Its role here is to compare the exponentially small physical
flux with its quadratic value \emph{relatively}, uniformly in the number
of reflections. The trace-norm bound in \eqref{eq:log-det} and the summable
endpoint bridge make this possible on one fixed neighborhood.

For a fixed $j$, ordinary Morse integration would yield a quadratic onset
with a constant depending on that segment. It would not, by itself, give
a common collar for all $j$, relative remainder bounds after division by
$\sinh(j\chi_R)^{-1}$, or uniform bounds for additive marks and sampling
cuts. Lemmas~\ref{lem:localization}, \ref{lem:bridge} and
\ref{lem:relative} give these additional statements for the complete
mechanical event. Theorem~\ref{thm:marked} then gives the summable
source-dependent hierarchy, and Section~\ref{sec:records} identifies
its joint record. The comparison with an independent roof having the same
marginal shows a change of onset exponent, not merely a different constant.

\subsection{Instability from statistical thresholds and from marked lengths}
B\'alint, De Simoi, Kaloshin and Leguil \cite{BDKL} recover curvatures at
period-two points and Lyapunov exponents of periodic orbits from the marked
length spectrum of open dispersing billiards with three strictly convex
obstacles satisfying the non-eclipse condition. Their datum consists of
labelled periodic lengths. Our coefficient sequence comes instead from
full-equilibrium probabilities of maximal collision counts in a periodic
circular billiard. Formula~\eqref{eq:ratio} recovers the normal
orbit's multiplier from these statistical thresholds. This is a different
observable, not a claim to stronger inverse-geometric rigidity. The
identification of the multiplier with the Jacobi recurrence is proved
explicitly in Section~\ref{sec:consequences}.


\subsection{Periodic rare events and physical collision onsets}
Carney, Nicol and Zhang \cite[Theorems 1--2]{CNZ} prove extreme-value
and compound-Poisson limits for visits to shrinking neighborhoods of
regular periodic points, including finite- and infinite-horizon Sinai
billiard maps. Their metric is adapted to stable and unstable foliations;
the extremal index is $1-|DT^p|_{E^u}|^{-1}$, and the associated
cluster multiplicities are geometric. Their observation length diverges
while the target shrinks so that the expected exceedance count stays
finite. Thus statistical detection of periodic instability is already
present in that work; it is not an originality claim of this article.

Our observable is the maximal number of physical impacts in a window
near $jg_e$, under full flow preparation. The collar, relative error and
shape/source derivatives are uniform in $j$, rather than derived from
a rare-visit limit. Theorems~\ref{thm:g-stability} and
\ref{thm:g-marked} verify those claims under noncircular deformations
and at competing threshold surfaces. Neither a rare-visit law nor a
fixed-segment Morse asymptotic supplies this uniform physical statement.
The proof does not invoke a mixing rate. Conversely it does not prove
a compound-Poisson law for long observation sequences.

The inverse issue is likewise explicit. In the circular model the
threshold spacing already fixes the radius and multiplier. We retain
Corollary~\ref{cor:multiplier} as its exact amplitude identity, not as
recovery of another unknown degree of freedom. In
Theorem~\ref{thm:g-inverse}, three independent curvature values vary
while every threshold location remains fixed. The unlabelled amplitudes
resolve precisely this uncertainty, with the unknown preparation area
removed by recovery from the same data. This does not assert full shape
rigidity or a stability estimate for noisy inverse data. It gives a
positive geometric consequence not determined by the threshold locations.

The marked-length results in \cite{BDKL} concern labelled periodic-orbit
lengths in open dispersing billiards and recover dynamical and curvature
information from those data. Our equal-gap family and coefficient
inversion distinguish the particular full-equilibrium data used here;
they do not purport to supersede that inverse theory. The variational
relation of action Hessians and multipliers remains classical
\cite{Hill}. The new relative estimate is its uniformly controlled use
for a complete physical threshold event with deformable geometry.

\subsection{Perturbative limit laws and the original Fourier programme}
Demers and Zhang \cite{DemersZhang} construct a functional-analytic
framework for billiard-map perturbations, including scatterer motions and
deformations. Their perturbation estimates and continuity of spectral
data address singularities on a much larger dynamical phase space than
our onset charts. Theorem~\ref{thm:marked} does not replace those operator
estimates. It gives every fixed radius derivative of a specified extremal
record, including its threshold singularity, in a regime selected by an
exact physical event. It is not an assertion of all-order differentiability
of the full transfer operator under arbitrary deformations.

Demers, Melbourne and Nicol \cite{DMN} prove martingale approximation and
statistical limit laws for H\"older observables of the time-one map of a
finite-horizon planar periodic Lorentz gas. Those typical-orbit limit laws
and the present maximal-count asymptotics concern different regimes.
In particular, neither their stated central limit and invariance principles
nor our threshold series is a source-differentiated mixed local Edgeworth
expansion for general collision records.

The earlier arithmetic and Fourier-inversion chapter \cite{A2R33} proves
an algebraic separation estimate and integrates four explicitly assumed
frequency envelopes. A mechanical application requires the characteristic
family itself to satisfy those envelopes, together with a central
integrable remainder and the required lattice and covariance conditions.
The present results provide a uniform, source-dependent calculation on a
nontrivial family of actual multi-impact events. They do not estimate the
complementary itineraries or construct the stable-curve operators needed
for that original long-time application. In particular, the sum in
\eqref{eq:summable} ranges over distinct onset windows; it cannot be
substituted for the characteristic function of one unrestricted long
record. Establishing that substitution's actual prerequisites is a
separate mathematical problem, not a consequence of the notation.

Thus the structural conclusion is a collision-order-uniform geometric
threshold law, including noncircular stability, competing onsets,
curvature identification and specified marked response under full preparation. The
fixed-window theorem in \cite{TC}, the conditional Fourier compiler in
\cite{A2R33}, and this uniform extremal theorem have separate, explicit
hypotheses and conclusions. No implication between them is used without
proof.

\begin{thebibliography}{9}
\bibitem{BDKL}
P.~B\'alint, J.~De Simoi, V.~Kaloshin and M.~Leguil,
Marked length spectrum, homoclinic orbits and the geometry of open
dispersing billiards, \emph{Comm. Math. Phys.} \textbf{374} (2020),
no.~3, 1531--1575.
\href{https://doi.org/10.1007/s00220-019-03448-x}{doi:10.1007/s00220-019-03448-x}.

\bibitem{Hill}
S.~V.~Bolotin and D.~V.~Treschev, Hill's formula,
\emph{Russian Math. Surveys} \textbf{65} (2010), no.~2, 191--257.
\href{https://doi.org/10.1070/RM2010v065n02ABEH004671}{doi:10.1070/RM2010v065n02ABEH004671}.

\bibitem{DMN}
M.~Demers, I.~Melbourne and M.~Nicol, Martingale approximations and
anisotropic Banach spaces with an application to the time-one map of a
Lorentz gas, \emph{Nonlinearity} \textbf{33} (2020), 4095--4113.
\href{https://doi.org/10.1088/1361-6544/ab7d22}{doi:10.1088/1361-6544/ab7d22}.

\bibitem{DemersZhang}
M.~F.~Demers and H.-K.~Zhang, A functional analytic approach to
perturbations of the Lorentz gas, \emph{Comm. Math. Phys.}
\textbf{324} (2013), no.~3, 767--830.
\href{https://doi.org/10.1007/s00220-013-1820-0}{doi:10.1007/s00220-013-1820-0}.

\bibitem{Marklof}
J.~Marklof, Entry and return times for semi-flows,
\emph{Nonlinearity} \textbf{30} (2017), no.~2, 810--824.
\href{https://doi.org/10.1088/1361-6544/aa518b}{doi:10.1088/1361-6544/aa518b}.

\bibitem{TC}
Q.~Qi, \emph{Two-collision counting response in a periodic Lorentz gas},
September 8, 2026, complete source-pinned companion to this article.
The reviewed original is retained without alteration as
\texttt{two\_collision.tex} and \texttt{two\_collision.pdf}.

\bibitem{A2R33}
Q.~Qi, \emph{A2: nonempty arithmetic and a checkable Fourier-inversion
budget}, Round 33 source reconstruction, September 3, 2026.
The original chapter and its source identity are retained in the
submission history; the conditional inversion hypotheses are not assumed
as established properties of the mechanical family in this article.
\bibitem{CNZ}
M. Carney, M. Nicol and H.-K. Zhang,
Compound Poisson law for hitting times to periodic orbits in two-dimensional
hyperbolic systems, \emph{J. Stat. Phys.} \textbf{169} (2017), 804--823.
\href{https://doi.org/10.1007/s10955-017-1893-9}{doi:10.1007/s10955-017-1893-9}.
Theorems 1--2 in the primary preprint, arXiv:1709.00530.
\end{thebibliography}

\end{document}
